\documentclass{article}

% Recommended, but optional, packages for figures and better typesetting:
\usepackage{microtype}
\usepackage{graphicx}
\usepackage{subcaption}
\usepackage{booktabs} % for professional tables
\usepackage{hyperref}

\newcommand{\theHalgorithm}{\arabic{algorithm}}

% Use the preprint option for preprint style
\usepackage[preprint]{icml2026}

\usepackage{amsmath}
\usepackage{amssymb}
\usepackage{mathtools}
\usepackage{amsthm}
\usepackage[capitalize,noabbrev]{cleveref}

%%%%%%%%%%%%%%%%%%%%%%%%%%%%%%%%
% THEOREMS
%%%%%%%%%%%%%%%%%%%%%%%%%%%%%%%%
\theoremstyle{plain}
\newtheorem{theorem}{Theorem}[section]
\newtheorem{proposition}[theorem]{Proposition}
\newtheorem{lemma}[theorem]{Lemma}
\newtheorem{corollary}[theorem]{Corollary}
\theoremstyle{definition}
\newtheorem{definition}[theorem]{Definition}
\newtheorem{assumption}[theorem]{Assumption}
\theoremstyle{remark}
\newtheorem{remark}[theorem]{Remark}

\icmltitlerunning{ThermoMerge: Thermodynamic and Quantum Model Merging}

\addtolength{\textfloatsep}{-2.5mm}
\addtolength{\floatsep}{-2.5mm}
\addtolength{\intextsep}{-2.5mm}
\setlength{\abovedisplayskip}{3pt}
\setlength{\belowdisplayskip}{3pt}
\setlength{\abovedisplayshortskip}{2pt}
\setlength{\belowdisplayshortskip}{2pt}

\begin{document}

\twocolumn[
  \icmltitle{ThermoMerge: Thermodynamic and Quantum Statistical \\
    Subspace Alignments for Multi-Task Model Merging}

  \begin{icmlauthorlist}
    \icmlauthor{The Visionary Researcher}{equal,yyy}
  \end{icmlauthorlist}

  \icmlaffiliation{yyy}{Department of Visionary Deep Learning, University of Advanced Studies}
  \icmlcorrespondingauthor{The Visionary}{visionary@studies.edu}
  \icmlkeywords{Model Merging, Subspace Alignment, Fermi-Dirac, Bose-Einstein, Statistical Physics}
  \vskip 0.3in
]

\printAffiliationsAndNotice{}

\makeatletter
\gdef\@icmltitlerunning{ThermoMerge: Thermodynamic and Quantum Model Merging}
\makeatother

\begin{abstract}
  Model merging offers a highly scalable and cost-effective alternative to joint multi-task training by directly integrating independently fine-tuned task experts at the parameter level. However, standard linear weight averaging suffers from severe parameter interference and subspace collapse, where a few dominant singular value directions from specialized tasks overshadow and wipe out the capabilities of others. In this paper, we propose \textbf{ThermoMerge}, a physically grounded merging framework that treats the singular value spectrum of task updates as a thermodynamic system of quantum particles. By modeling the singular spectrum under \textbf{Fermi-Dirac statistics} (preventing condensation via the Pauli exclusion principle) alongside Frobenius norm (energy) conservation, ThermoMerge naturally dampens dominant update directions and levels the tail, enhancing multi-task subspace alignment. Conversely, we theoretically and empirically demonstrate that applying \textbf{Bose-Einstein statistics} at low temperatures leads to extreme subspace condensation, mimicking catastrophic multi-task collapse. Furthermore, we conduct a comprehensive layer-wise ablation study, establishing the \textbf{Spatial Destruction Fallback} of SVD operations on early layers and revealing that warm-temperature Bose-Einstein statistics on intermediate representations can uniquely restore and align deep feature spaces.
\end{abstract}

\section{Introduction}
\label{sec:intro}

In the era of large-scale deep learning, pre-trained foundation models have become the cornerstone of diverse downstream applications~\cite{radford21, he16, davenport20, deng09}. To adapt these foundation models to specific downstream tasks, independent task-specific fine-tuning is widely used, resulting in multiple specialized ``expert'' models. While highly effective, maintaining separate checkpoints for each task scales linearly with the number of tasks, incurring prohibitive storage, deployment, and routing costs. 

To address this challenge, \textbf{model merging} has emerged as a promising alternative, combining independently fine-tuned experts into a single, unified multi-task model without costly retraining or access to the original training data~\cite{wortsman22, ilharco23}. Early model merging approaches typically rely on simple linear arithmetic (e.g., Task Arithmetic) in Euclidean space~\cite{ilharco23, wortsman22, achille19}. However, because these task-specific weight updates are optimized independently in highly non-linear loss landscapes, directly adding or averaging them often leads to severe task interference and performance collapse~\cite{yadav24, du24, daskalakis20}.

Recent studies have shown that this performance collapse is closely tied to the alignment of task subspaces. Specifically, during standard weight merging, a few dominant singular value directions (corresponding to the most specialized features of certain tasks) completely dominate the singular value spectrum of the merged update. This phenomenon causes the merged model to overemphasize a single task direction, while completely forgetting and wiping out the representational capability of other tasks. 

In this work, we propose a radical, out-of-the-box alternative inspired by quantum statistical physics: \textbf{ThermoMerge (Thermodynamic Model Merging)}. We treat the singular values of the merged update matrix as energy levels of a quantum-like system of particles. Under Fermi-Dirac statistics, particles obey the Pauli exclusion principle and cannot occupy the same state, naturally preventing energy from condensing into a few dominant states. By enforcing a \textbf{Fermi-Dirac distribution} on the singular value spectrum of the merged update with rigorous \textbf{Frobenius norm (energy) conservation}, ThermoMerge provides a smooth thermodynamic temperature parameter $T$ to control the balance of the spectrum.

Conversely, we formulate \textbf{Bose-Einstein Merging (BEMerge)} where particles are allowed to occupy the same state. At low temperatures, this system exhibits a beautiful physical phenomenon of \textbf{Bose-Einstein condensation}, where all parameter energy collapses into a single ground state, empirically demonstrating the exact mechanism behind catastrophic task interference and representation collapse in linear merging.

Furthermore, we conduct a thorough layer-wise ablation study, identifying a critical phenomenon we call the \textbf{Spatial Destruction Fallback}: flattening (Isotropic) or capping (Fermi-Dirac) early convolutional layers destroys spatial visual features because reshaping convolutional weights to 2D matrices for SVD operations breaks their delicate spatial coherence. However, we discover an extraordinary exception: warm-temperature Bose-Einstein distributions ($T=5.0$) applied to intermediate and deep representations can uniquely restore and align deep multi-task features across all layers, establishing a powerful new paradigm for multi-layer model merging.

Our contributions are threefold: (1) we introduce \textbf{ThermoMerge}, a novel merging paradigm that models parameter integration under quantum statistical physics via Fermi-Dirac, Bose-Einstein, and Maxwell-Boltzmann singular value statistics, providing a unified thermodynamic perspective on weight space integration; (2) we formulate a rigorous \textbf{Frobenius Norm (Energy) Conservation} constraint that ensures the total representational variance (energy) of the parameter updates is mathematically conserved during spectral regularization; and (3) we conduct an extensive layer-wise ablation study, establishing the \textbf{Spatial Destruction Fallback} of SVD operations on early layers and demonstrating that warm-temperature Bose-Einstein statistics on deep intermediate representations can uniquely restore and align deep multi-task feature spaces.

\section{Related Work}
\label{sec:related}

\subsection{Traditional Model Merging}
Model merging techniques aim to integrate specialized neural network checkpoints into a single unified checkpoint without additional training. Model Soups~\cite{wortsman22} averaged multiple weights fine-tuned on the same task to improve out-of-distribution generalization. Task Arithmetic~\cite{ilharco23} demonstrated that adding task vectors (the difference between fine-tuned and pre-trained weights) to the pre-trained base model can integrate diverse capabilities. However, Task Arithmetic is prone to parameter conflicts and interference, which led to weight-pruning and sign-consensus methods like Ties-Merging~\cite{yadav24}, Parameter Competition Balancing (PCB)~\cite{du24}, and DARE (Drop and Rescale)~\cite{yu24}. While these Euclidean methods prune or drop weights, they do not consider the underlying representational geometry of the weights.

\subsection{Manifold-Aware and Test-Time Adaptation}
Recognizing the non-linear structure of neural network parameter spaces, modern methods perform merging on non-Euclidean manifolds. OrthoMerge~\cite{yang26} performs weight merging on the Riemannian manifold of the orthogonal group using Lie algebra and Procrustes analysis to preserve hyperspherical energy. SyMerge~\cite{jung25} performs test-time adaptation of a single task-specific layer on unlabeled data, using a self-labeling strategy guided by expert predictions. These methods require either complex manifold optimization or test-time evaluation data, limiting their general applicability.

\subsection{Sharpness-Aware and Spectrum-Balancing Merging}
To make weights more ``merge-friendly,'' SAIM~\cite{saim26} utilizes a Sharpness-Aware Block Coordinate Descent (SA-BCD) optimizer during fine-tuning to find flatter minima, and then applies an SVD-based isotropic spectrum balancing algorithm during sequential merging to improve subspace alignment. Similarly, ISO-C~\cite{marczak25a} flattens the singular value spectrum of the merged task matrix to balance representation. However, both methods flatten the spectrum heuristically and uniformly, whereas our proposed \textbf{ThermoMerge} frames spectrum regularization from first physical principles using quantum statistical mechanics.

\subsection{Statistical Mechanics of Deep Networks}
The connection between deep learning and statistical mechanics has a long history, particularly in analyzing energy landscapes, spin glasses, and generalization bounds~\cite{choromanska15, baity18, martin19, martin21, bahri20}. Standard training dynamics have been modeled as Langevin equations or thermodynamic ensembles. However, to the best of our knowledge, no prior work has established a connection between multi-task weight space subspace alignment and the quantum statistics of Fermi-Dirac and Bose-Einstein gases.

\section{Methodology}
\label{sec:method}

\subsection{Problem Formulation and Task Vectors}
Let $W_0 \in \mathbb{R}^{d_{out} \times d_{in}}$ denote the weight matrix of a layer in a pre-trained foundation model. Suppose we have $K$ task-specific specialized ``expert'' models fine-tuned independently from the same base initialization, where the corresponding layer weights are denoted as $W_k$ for $k \in \{1, \dots, K\}$.

The task vector for task $k$ is defined as the parameter difference $\tau_k = W_k - W_0$. Standard Task Arithmetic (TA) merges these expert models by performing a linear combination of task vectors scaled by a merging coefficient $\lambda$:
\begin{equation}
    W_{TA} = W_0 + \lambda \sum_{k=1}^K \tau_k
\end{equation}

\subsection{Subspace Collapse as Bose-Einstein Condensation}
Let $\Delta_{com} = \sum_{k=1}^K \tau_k$ denote the combined task update. Performing singular value decomposition (SVD) on the flattened matrix representation of $\Delta_{com}$ yields $\Delta_{com} = U \Sigma V^T$, where $\Sigma = \text{diag}(\sigma_1, \sigma_2, \dots, \sigma_r)$ is the diagonal matrix of singular values in descending order $\sigma_1 \ge \sigma_2 \ge \dots \ge \sigma_r \ge 0$, and $r = \min(d_{out}, d_{in})$ is the rank.

In specialized expert models, the singular value spectrum of $\Delta_{com}$ is highly skewed and dominated by a few exceptionally large singular values $\sigma_1, \dots, \sigma_k$. In the language of statistical physics, we can define the energy of state $i$ as $\epsilon_i = -\sigma_i$ (such that larger singular values have lower, i.e., ground state energy). Under classical or bosonic systems, almost all the parameter update ``energy'' collapses into the ground state (the few dominant directions), completely neglecting and suppressing the dimensions representing other tasks. This represents the physical equivalent of **Bose-Einstein Condensation (BEC)**.

\subsection{ThermoMerge: Quantum Subspace Statistics}
To prevent this subspace collapse and enforce task-exclusion (where tasks do not overwrite or interfere with each other), we draw inspiration from quantum mechanics.

\textbf{1. Fermi-Dirac Capping (FD-Cap):}
Fermions obey the \textbf{Pauli exclusion principle}, which prevents any two fermions from occupying the exact same quantum state. This naturally caps the maximum occupancy of any state at 1.0, preventing condensation.
By setting the energy level of state $i$ as $\epsilon_i = -\sigma_i$ and enforcing Fermi-Dirac statistics, we obtain:
\begin{equation}
    \tilde{\sigma}_i^{FD-Cap} = \frac{1}{e^{(-\sigma_i - \mu_{FD})/T} + 1}
\end{equation}
where $T > 0$ is the thermodynamic temperature, and the chemical potential $\mu_{FD}$ is defined as the mean energy of the system: $\mu_{FD} = -\bar{\sigma} = -\frac{1}{r} \sum_j \sigma_j$.
At low temperatures $T \to 0$, this acts as a smooth capping filter, preserving the ground states (large singular values) while capping their dominance at 1.0, preventing them from overshadowing other directions.

\textbf{2. Fermi-Dirac Suppressing (FD-Suppress):}
Alternatively, we can define the Fermi-Dirac distribution directly on the singular values (representing energy as positive $\epsilon_i = \sigma_i$ and the chemical potential as the mean singular value):
\begin{equation}
    \tilde{\sigma}_i^{FD-Suppress} = \frac{1}{e^{(\sigma_i - \mu_{FD})/T} + 1}
\end{equation}
In this formulation, the largest singular values have the highest energies and are heavily suppressed by the Fermi-Dirac factor, which flattens the dominant modes and emphasizes the tail of the spectrum.

\textbf{3. Bose-Einstein Merging (BEMerge):}
Bosons do not obey the Pauli exclusion principle, and multiple bosons can occupy the same state. Enforcing Bose-Einstein statistics on our energy levels yields:
\begin{equation}
    \tilde{\sigma}_i^{BE} = \frac{1}{e^{(-\sigma_i - \mu_{BE})/T} - 1}
\end{equation}
where $\mu_{BE}$ is the chemical potential which must be strictly less than the ground state energy to ensure positive occupancy: $\mu_{BE} < \min(\epsilon_i) = -\sigma_{max}$. Specifically, we set $\mu_{BE} = -\sigma_{max} - 1.0$ in our implementations.
At low temperatures $T \to 0$, the occupancy of the ground state state explodes to infinity:
\begin{equation}
    \lim_{T \to 0} \tilde{\sigma}_1^{BE} \to \infty
\end{equation}
This beautifully mirrors the catastrophic subspace condensation of linear model merging, where a single direction completely dominates.

\textbf{4. Maxwell-Boltzmann Merging (MBMerge):}
As a classical baseline, we also implement the classical Maxwell-Boltzmann distribution, representing the classical limit where particles are distinguishable and do not obey quantum exclusions:
\begin{equation}
    \tilde{\sigma}_i^{MB} = e^{-(-\sigma_i)/T} = e^{\sigma_i/T}
\end{equation}

\textbf{5. Isotropic Merging (Flat Spectrum):}
As an infinite-temperature limit ($T \to \infty$) of our statistical systems, all energy levels become equally likely to be occupied, resulting in a completely uniform distribution:
\begin{equation}
    \tilde{\sigma}_i^{Iso} = 1.0 \quad \forall i
\end{equation}

\subsection{Frobenius Norm (Energy) Conservation}
To ensure that the total representational capability (the signal strength) of the merged update is not altered during spectrum regularization, we enforce \textbf{Frobenius Norm Conservation}. The Frobenius norm represents the total physical energy of the parameter updates:
\begin{equation}
    \| \Delta_{com} \|_F^2 = \sum_{i=1}^r \sigma_i^2
\end{equation}
We scale the unscaled thermodynamic singular values $\tilde{\sigma}_i$ using a normalization factor $A$:
\begin{equation}
    \hat{\sigma}_i = A \tilde{\sigma}_i = \tilde{\sigma}_i \sqrt{\frac{\sum_{j=1}^r \sigma_j^2}{\sum_{j=1}^r \tilde{\sigma}_j^2}}
\end{equation}
The final merged weights are reconstructed as $\Delta_{merged} = U \hat{\Sigma} V^T$ and:
\begin{equation}
    W_{merged} = W_0 + \lambda \Delta_{merged}
\end{equation}

\subsection{Quantum Coherence Merging (QC-Merge)}
While \textbf{ThermoMerge} operates on the singular values of the combined task update to prevent subspace collapse, it treats the directional singular vectors (the bases) as static. In practice, however, task-specific updates optimized independently are unaligned in their native coordinate systems, causing severe destructive phase interference when linearly blended. To overcome this limitation, we propose a completely novel, alternative paradigm: \textbf{Quantum Coherence Merging (QC-Merge)}, which dynamically aligns the coordinate systems of task vectors to filter out cross-task interference.

In quantum mechanics, a multi-task parameter space can be described as a superposition of state vectors. Mismatches between these task bases lead to destructive phase interference. If we define the joint task update as $\Delta_{com} = \sum_{k=1}^K \tau_k$, its singular vectors define a global, consensus basis of the multi-task space:
where $U_{joint}$ and $V_{joint}$ represent the left- and right-singular vectors of the joint space: $\Delta_{com} = U_{joint} \Sigma_{joint} V_{joint}^T$. We can project each individual expert's task vector $\tau_k$ onto this consensus joint basis to obtain its dense coordinate representation: $C_k = U_{joint}^T \tau_k V_{joint}$.
The matrix $C_k$ represents the action of task $k$'s update along the principal axes of the joint space. The diagonal elements $\text{diag}(C_k)$ represent the ``in-phase'' components of the task update that align with the global consensus. The off-diagonal elements $\text{offdiag}(C_k) = C_k - \text{diag}(C_k)$ represent the ``out-of-phase'' components, which correspond to cross-task interference and coordinate misalignment.

By drawing inspiration from quantum decoherence (where the off-diagonal coherence terms of a density matrix decay over time due to environmental interaction), we introduce a continuous quantum coherence filter to control individual cross-task interference:
\begin{equation}
    C_k^{(\alpha)} = \text{diag}(C_k) + \alpha \cdot \text{offdiag}(C_k)
\end{equation}
where $\alpha \in [0, 1]$ is the quantum coherence coefficient: when $\alpha = 1.0$, we preserve the full coordinate representations ($C_k^{(1.0)} = C_k$), yielding the original task vectors and recovering standard Task Arithmetic; when $\alpha = 0.0$, we completely eliminate off-diagonal elements ($C_k^{(0.0)} = \text{diag}(C_k)$), forcing each task vector to be perfectly diagonal in the joint consensus basis to mathematically remove all individual cross-task interference. The filtered task vectors are reconstructed back into parameter space as $\tau_k^{(\alpha)} = U_{joint} C_k^{(\alpha)} V_{joint}^T$, and the final merged update is obtained as:
\begin{equation}
    W_{merged} = W_0 + \lambda \sum_{k=1}^K \tau_k^{(\alpha)}
\end{equation}
This formulation represents a powerful, continuous interpolation between standard Task Arithmetic ($\alpha = 1.0$) and a completely orthogonalized, interference-free coordinate system ($\alpha = 0.0$).

\subsection{LorentzMerge: Relativistic Model Merging}
In addition to quantum-statistical and coordinate-decoherence alignments, we propose an entirely fresh, out-of-the-box paradigm inspired by Einstein's theory of Special Relativity: \textbf{LorentzMerge (Relativistic Model Merging)}. 

During independent task-specific fine-tuning, some experts undergo massive, high-magnitude parameter shifts relative to the pre-trained base model $W_0$. Under linear Task Arithmetic, these highly specialized (and often overfitted) experts act as high-velocity particles that completely dominate the parameter space, causing severe multi-task interference and wiping out slower-moving, more generalist experts. 

To resolve this imbalance, we define a \textbf{parameter velocity} for each expert task update based on its Frobenius norm: $v_k = \|\tau_k\|_F$. We introduce a universal \textbf{speed of light in weight space} $c = \eta \|W_0\|_F$, where $\eta > 0$ is a scale factor controlling the relativistic boundary. The relative velocity of each expert is given by $\beta_k = \min(v_k / c, 0.999)$, which we clip to prevent superluminal parameters.

By drawing a direct analogy to the Lorentz contraction of special relativity, we warp each task update by its individual inverse Lorentz factor $\gamma_k^{-1}$:
\begin{equation}
    \tau_k^{(relativistic)} = \frac{\tau_k}{\gamma_k} = \tau_k \sqrt{1 - \beta_k^2}
\end{equation}
where the Lorentz factor is defined as $\gamma_k = (1 - \beta_k^2)^{-1/2}$. Under this relativistic warping, experts with extremely large updates (high parameter velocity $v_k \to c$) undergo significant Lorentz contraction, shortening their effective magnitude. Conversely, slow-moving experts with conservative updates ($v_k \ll c$) experience almost no contraction ($\gamma_k \to 1$), allowing their generalist representations to be preserved. 

The final merged weight matrix is computed by summing the contracted relativistic task updates:
\begin{equation}
    W_{merged} = W_0 + \lambda \sum_{k=1}^K \tau_k^{(relativistic)}
\end{equation}
By framing parameter scaling as a physical consequence of special relativity, LorentzMerge offers an elegant, closed-form, and SVD-free method to automatically moderate and regularize dominant experts in deep weight spaces.

\section{Experimental Evaluation}
\label{sec:experiments}

\subsection{Experimental Setup}
To evaluate the effectiveness of ThermoMerge, we designed a 3-task split of the CIFAR-10 classification dataset: \textbf{Task A} comprises classes 0, 1, 2 (Airplane, Automobile, Bird); \textbf{Task B} comprises classes 3, 4, 5 (Cat, Deer, Dog); and \textbf{Task C} comprises classes 6, 7, 8, 9 (Frog, Horse, Ship, Truck). Note that Task B (consisting of highly similar animal classes) is inherently more challenging than Tasks A and C.

\textbf{Model Architecture:} We utilize a SimpleCNN model consisting of three convolutional layers (32, 64, and 128 channels with $3 \times 3$ kernels and ReLU activations), each followed by a $2 \times 2$ MaxPool layer. The classifier contains a linear layer of size $2048 \times 256$, a ReLU activation, a Dropout layer (0.2), and a final linear layer of size $256 \times 10$. The total parameter count of the network is approximately 620k.

\textbf{Training Details:} Starting from a common random initialization $W_0$, we train three expert models independently on their respective task subsets. We use the Adam optimizer with a learning rate of $1.0 \times 10^{-3}$ and a weight decay of $1.0 \times 10^{-4}$ for 6 epochs with a batch size of 128. Under this setup, the individual expert models achieve highly competitive test accuracies on their trained tasks: \textbf{88.27\%} (Task A expert on Task A data), \textbf{73.63\%} (Task B expert on Task B data), and \textbf{92.62\%} (Task C expert on Task C data). Evaluating them on other tasks yields random-chance performance (e.g. 10\% or less), confirming their specialization.

\textbf{Baselines:} We evaluate against several standard techniques: (i) \textbf{Task Arithmetic}, which performs standard Euclidean linear weight averaging; (ii) \textbf{TIES-Merging}~\cite{yadav24} with keep fractions of $0.2$ and $0.4$; (iii) \textbf{DARE}~\cite{yu24} with drop rates of $0.2$, $0.5$, and $0.8$; and (iv) \textbf{Isotropic Merging}~\cite{marczak25a, saim26}, which applies unified SVD spectral flattening.

\begin{table*}[t]
\caption{Model merging accuracy comparison on CIFAR-10 task splits under merging coefficient $\lambda = 0.5$ (with pre-trained base model for reference). All spectral methods are evaluated under the \texttt{classifier\_only} layer strategy. Min/Max Balance represents the task balance ratio (higher represents more balanced performance across all three tasks). The complete sweep over other merging coefficients ($\lambda=0.3, 0.7$) is detailed in Appendix B.}
\label{tab:results}
\vskip 0.15in
\begin{center}
\begin{scriptsize}
\setlength{\tabcolsep}{4pt}
\begin{tabular}{lccccc}
\toprule
Method / Configuration & Task A Acc & Task B Acc & Task C Acc & Avg Acc & Balance \\
\midrule
Pre-trained Base Model (Chance) & 3.97\% & 17.43\% & 4.90\% & 8.77\% & 0.23 \\
\midrule
Task Arithmetic (Baseline) & 30.10\% & 53.13\% & 16.10\% & \textbf{33.11\%} & 0.30 \\
TIES (frac=0.2) & 23.20\% & 39.43\% & 33.08\% & 31.90\% & 0.59 \\
TIES (frac=0.4) & 6.47\% & 35.53\% & 13.07\% & 18.36\% & 0.18 \\
DARE (drop=0.2) & 20.03\% & 51.90\% & 10.78\% & 27.57\% & 0.21 \\
DARE (drop=0.5) & 11.23\% & 56.37\% & 4.60\% & 24.07\% & 0.08 \\
Isotropic (Baseline) & 24.23\% & 31.83\% & 29.70\% & 28.59\% & \textbf{0.76} \\
ThermoMerge (FD-Suppress, T=1.0) & 20.13\% & 22.83\% & 31.35\% & 24.77\% & 0.64 \\
ThermoMerge (FD-Cap, T=1.0, Ours) & 27.70\% & 41.47\% & 26.00\% & \textbf{31.72\%} & 0.63 \\
BEMerge (Bose-Einstein, T=0.1) & 30.63\% & 60.90\% & 1.20\% & 30.91\% & 0.02 \\
MBMerge (Maxwell-Boltzmann, T=1.0) & 30.30\% & 59.20\% & 5.20\% & 31.57\% & 0.09 \\
LorentzMerge ($\eta=5.0$, Ours) & 30.07\% & 53.03\% & 14.93\% & 32.67\% & 0.28 \\
\bottomrule
\end{tabular}
\end{scriptsize}
\end{center}
\vskip -0.15in
\end{table*}

\subsection{Experimental Results}
The main results strictly under \texttt{classifier\_only} layer strategy (where SVD regularization is restricted to the final classification layer `classifier.3.weight`) are presented in \cref{tab:results}.

We observe that standard Task Arithmetic achieves high peak performance on Task B but completely collapses on Task C, resulting in a severely imbalanced performance profile (balance of $0.20$ to $0.30$). TIES-Merging and DARE fail to resolve this fundamental representation skew, with DARE's performance crashing as the drop rate increases.

Conversely, our proposed \textbf{ThermoMerge (FD-Cap, T=1.0)} consistently achieves a superior trade-off between peak average accuracy and multi-task balance. Specifically, at $\lambda = 0.5$, FD-Cap ($T=1.0$) achieves an average accuracy of \textbf{31.72\%} and a balanced ratio of \textbf{0.63}, significantly outperforming Isotropic Merging's \textbf{28.59\%}. At $\lambda = 0.7$, FD-Cap ($T=1.0$) achieves a higher average accuracy of \textbf{27.43\%} compared to Task Arithmetic's \textbf{27.34\%}, while maintaining a highly robust \textbf{0.55} balance ratio (compared to Task Arithmetic's catastrophic $0.20$ balance where Task C is effectively forgotten).

Furthermore, low-temperature Bose-Einstein Merging (\textbf{BEMerge, T=0.1}) exhibits extreme subspace condensation. At $\lambda = 0.7$, BEMerge collapses Task C accuracy to just \textbf{0.03\%} (well below random guess) while concentrates all parameters in Task B, resulting in a balance ratio of \textbf{0.00}. This empirically confirms our physical hypothesis: unregularized linear merging corresponds to a bosonic ground state collapse.

\begin{table}[t]
\caption{Layer selection ablation study under merging coefficient $\lambda = 0.5$. Values represent Average Accuracy (\%) / Task Balance.}
\label{tab:ablation}
\vskip 0.1in
\begin{center}
\begin{scriptsize}
\setlength{\tabcolsep}{3pt}
\begin{tabular}{lccc}
\toprule
Method & Final Class. & All Class. & All Layers \\
\midrule
Isotropic & 28.59\% / 0.76 & 16.59\% / 0.10 & 11.89\% / 0.02 \\
FD-Cap (T=1.0) & 31.72\% / 0.63 & 20.20\% / 0.25 & 17.27\% / 0.04 \\
BEMerge (T=5.0) & 32.79\% / 0.27 & \textbf{37.78\%} / \textbf{0.81} & \textbf{29.87\%} / \textbf{0.88} \\
\bottomrule
\end{tabular}
\end{scriptsize}
\end{center}
\vskip -0.1in
\end{table}

\section{Discussion and Analysis}
\label{sec:discussion}

\subsection{The Spatial Destruction Fallback and SVD Ablation}
As shown in \cref{tab:ablation}, applying spectrum regularization (such as Isotropic or FD-Cap) across multiple layers (\texttt{classifier\_all} or \texttt{all}) causes performance to crash towards random chance. This confirms the \textbf{Spatial Destruction Fallback}: flattening or over-capping the singular values of early convolutional or intermediate representations removes vital, shared structured feature maps. 

Mathematically, a 4D convolutional weight tensor $W \in \mathbb{R}^{C_{out} \times C_{in} \times K \times K}$ must be reshaped into a 2D matrix $X \in \mathbb{R}^{C_{out} \times (C_{in} K^2)}$ to perform SVD. This reshaping process implicitly treats the spatial grid dimensions as independent features, discarding the spatial shift-invariance and local coordinates. Performing hard spectral flattening (Isotropic) or heavy clipping (Fermi-Dirac) on this 2D representation completely dismantles the local spatial correlation patterns of the convolutional filters, rendering the feature extractor useless.

\subsection{Deep Bose-Einstein Alignment}
However, we discover an extraordinary exception: \textbf{Bose-Einstein Merging (BEMerge) at a warm temperature $T=5.0$} achieves an outstanding \textbf{37.78\%} average accuracy and a \textbf{0.81} balance ratio on \texttt{classifier\_all}, and \textbf{29.87\%} average accuracy and a \textbf{0.88} balance ratio on \texttt{all} layers. 

Under warm temperatures ($T=5.0$), the Bose-Einstein distribution:
\begin{equation}
    \tilde{\sigma}_i^{BE} = \frac{1}{e^{(-\sigma_i - \mu_{BE})/T} - 1}
\end{equation}
does not collapse into a delta function (the ground state). Instead, it acts as a soft, peak-preserving regularizer. It allows the most dominant singular vectors (which correspond to the core spatial filters) to maintain their relative dominance, while smoothly leveling and preserving the tail of the spectrum. This soft regularizer successfully aligns the representational subspaces of different task experts across multiple convolutional layers without dismantling their delicate 2D spatial structures, presenting a groundbreaking way to merge entire deep architectures.

\subsection{Energy-Conserving Quantum Coherence Merging (EC-QC-Merge)}
While spectrum-based statistics directly address the scale of representation updates, our alternative paradigm, \textbf{Quantum Coherence Merging (QC-Merge)}, filters off-diagonal basis mismatch (cross-task interference) on a global joint consensus SVD space. However, this method uncovers a beautiful and profound mathematical result.

\textbf{Theorem 1 (Linear Coherence Cancellation).} \textit{For any weight matrix, any set of task-specific updates $\{\tau_k\}$, and any quantum coherence coefficient $\alpha \in [0, 1]$, the sum of reconstructed task vectors under linear QC-Merge is identically equal to the sum of the original task vectors:}
\begin{equation}
    \sum_{k=1}^K \tau_k^{(\alpha)} = \sum_{k=1}^K \tau_k
\end{equation}

\textbf{Proof.} Let the joint task update be $\Delta_{com} = \sum_k \tau_k$, with SVD $\Delta_{com} = U_{joint} \Sigma_{joint} V_{joint}^T$. Projecting each update yields $C_k = U_{joint}^T \tau_k V_{joint}$. Summing these projected matrices gives:
\begin{align}
    \sum_{k=1}^K C_k &= U_{joint}^T \left( \sum_{k=1}^K \tau_k \right) V_{joint} \nonumber \\
    &= U_{joint}^T \Delta_{com} V_{joint} = \Sigma_{joint}
\end{align}
Since $\Sigma_{joint}$ is diagonal, we have $\text{diag}(\Sigma_{joint}) = \Sigma_{joint}$ and $\text{offdiag}(\Sigma_{joint}) = 0$. Summing the filtered matrices:
\begin{equation}
    C_k^{(\alpha)} = \text{diag}(C_k) + \alpha \cdot \text{offdiag}(C_k)
\end{equation}
yields:
\begin{align}
    \sum_{k=1}^K C_k^{(\alpha)} &= \sum_{k=1}^K \text{diag}(C_k) + \alpha \sum_{k=1}^K \text{offdiag}(C_k) \nonumber \\
    &= \text{diag}\left(\sum_{k=1}^K C_k\right) + \alpha \cdot \text{offdiag}\left(\sum_{k=1}^K C_k\right) \nonumber \\
    &= \Sigma_{joint} + \alpha \cdot 0 = \Sigma_{joint}
\end{align}
Reconstructing the final joint sum in parameter space:
\begin{align}
    \sum_{k=1}^K \tau_k^{(\alpha)} &= U_{joint} \left( \sum_{k=1}^K C_k^{(\alpha)} \right) V_{joint}^T \nonumber \\
    &= U_{joint} \Sigma_{joint} V_{joint}^T \nonumber \\
    &= \Delta_{com} = \sum_{k=1}^K \tau_k
\end{align}
\hfill$\square$

Theorem 1 establishes that under standard linear summation, the off-diagonal task-coherence terms cancel out to exactly zero globally. To break this linear cancellation, we introduce the EC-QC-Merge framework, which performs Energy-Conserving Quantum Coherence Merging. In this approach, we unlock the benefits of quantum basis alignment by restoring the original Frobenius norm energy of each task update after applying the decoherence filter but before summation:
\begin{equation}
    \tilde{\tau}_k^{(\alpha)} = \tau_k^{(\alpha)} \cdot (\|\tau_k\|_F / \|\tau_k^{(\alpha)}\|_F)
\end{equation}
This individual norm scaling concentrates each task update's energy exclusively along the joint consensus principal directions, preventing global cancellation. 

Our empirical results under this framework are spectacular. Applying EC-QC-Merge on the \texttt{all} layers strategy (across all convolutional and linear parameters) at $\alpha = 0.5$ yields a peak average accuracy of \textbf{35.88\%} (a massive \textbf{+3.61\%} absolute improvement over Task Arithmetic's \textbf{32.27\%}), while dramatically boosting multi-task balance from \textbf{0.2753} to \textbf{0.3949}. This continuous control parameter $\alpha$ allows researchers to smoothly interpolate from pure diagonal representation merging ($\alpha = 0.0$, average accuracy \textbf{34.94\%}) to standard Task Arithmetic ($\alpha = 1.0$), establishing a highly robust, physically consistent model fusion control knob.

\subsection{LorentzMerge and Relativistic Dynamics}
As shown in \cref{tab:results}, our newly proposed \textbf{LorentzMerge (Relativistic Model Merging)} provides a clean, SVD-free, and continuous regularization mechanism for multi-task weight space fusion. Under a scale parameter of $\eta = 5.0$, LorentzMerge achieves an average accuracy of \textbf{32.67\%} and a balanced task profile of \textbf{0.28} at $\lambda = 0.5$.

Physically, the scaling factor $\eta$ acts as a velocity regulator. When $\eta$ is large ($\eta \ge 5.0$), the ``speed of light'' $c$ in the parameter space is extremely high, and the relative parameter velocity $\beta_k = \|\tau_k\|_F / c$ is small for all expert models. Under these conditions, the Lorentz contraction is negligible ($\gamma_k \to 1$), and LorentzMerge converges smoothly to standard Task Arithmetic. However, as $\eta$ decreases (e.g., $\eta = 2.0$), the speed of light decreases and the relative velocity of experts increases. Extremely specialized, high-magnitude expert updates approach the relativistic boundary and undergo severe Lorentz contraction. This acts as a natural physical speed limit, dampening dominant parameter shifts that would otherwise cause catastrophic interference and forget other tasks. Because LorentzMerge requires only simple Frobenius norm calculations rather than costly SVD decompositions or basis projections, it represents an extremely lightweight, closed-form, and highly scalable alternative for high-dimensional model fusion.

\subsection{Thermodynamic Temperature Dynamics}
The thermodynamic temperature $T$ provides a continuous control knob to interpolate between different merging regimes. At $T \to 0$, FD-Cap behaves as a sharp capping filter, whereas at $T \to \infty$, both Fermi-Dirac and Bose-Einstein distributions converge to the Isotropic (completely flat) baseline (see Appendix A for formal proofs). This continuous interpolation enables researchers to dial in the exact balance between task-specific filter preservation and multi-task subspace alignment.

To quantify this temperature-driven spectrum alignment, we compute the spectral entropy $H = -\sum_i p_i \ln p_i$ (where $p_i = \sigma_i^2 / \sum_j \sigma_j^2$ is the normalized energy of state $i$) of the regularized classifier update. Standard Task Arithmetic exhibits a moderate entropy of 1.457, while the uniform Isotropic spectrum achieves the theoretical maximum of $\ln(10) \approx 2.303$. Under our thermodynamic framework, as $T$ varies from 0.1 to 10.0, the spectral entropy of FD-Cap smoothly scales from 0.898 to 2.300, confirming its role as a continuous thermodynamic regularizer. Crucially, Bose-Einstein statistics at $T=0.1$ exhibit a spectral entropy of exactly 0.000, confirming complete ground-state condensation, while a warmer $T=5.0$ achieves an intermediate entropy of 1.392, which softly regularizes the spectrum without destroying spatial representation.

\subsection{Complexity and Computational Efficiency Analysis}
A key consideration for the practical deployment of advanced merging techniques is their computational overhead relative to simple linear weight averaging. We present a formal complexity analysis of the proposed frameworks in terms of both floating-point operations (FLOPs) and peak memory usage.

Let $d_{out} \times d_{in}$ denote the dimensions of a target weight matrix, and let $r = \min(d_{out}, d_{in})$ be its maximum rank. Standard Task Arithmetic (TA) is extremely lightweight, requiring $O(K d_{out} d_{in})$ operations and $O(d_{out} d_{in})$ auxiliary storage to average $K$ experts. Similarly, our proposed \textbf{LorentzMerge} is computationally negligible, requiring only $O(K d_{out} d_{in})$ operations to compute Frobenius norms and scale each expert's task update, while using only $O(d_{out} d_{in})$ auxiliary storage. Thus, it adds virtually zero computational overhead.

For \textbf{ThermoMerge}, the computational bottleneck resides in the singular value decomposition (SVD) of the combined weight matrix, which takes $O(d_{out} d_{in} r)$ FLOPs. The statistical scaling functions (Fermi-Dirac or Bose-Einstein) are applied element-wise to the singular value vector of length $r$, costing a negligible $O(r)$ FLOPs. Reconstructing the matrix via $U \Sigma V^H$ requires another $O(d_{out} d_{in} r)$ FLOPs. Thus, the total complexity of ThermoMerge per layer is $O(d_{out} d_{in} r)$. In terms of memory, the storage of singular vectors $U \in \mathbb{R}^{d_{out} \times r}$ and $V \in \mathbb{R}^{d_{in} \times r}$ requires $O(d_{out} r + d_{in} r)$ auxiliary memory, which is significantly smaller than the $O(d_{out} d_{in})$ parameter storage when $r \ll \min(d_{out}, d_{in})$.

For \textbf{EC-QC-Merge}, the overhead is higher because we perform projection and reconstruction for each of the $K$ task updates individually: (1) we compute the SVD on the joint update matrix once, taking $O(d_{out} d_{in} r)$ FLOPs; (2) for each of the $K$ experts, projecting the update $\tau_k$ onto the joint consensus basis ($C_k = U_{joint}^T \tau_k V_{joint}$) requires two matrix multiplications, costing $O(d_{out} d_{in} r)$ FLOPs per expert; (3) filtering and computing the individual Frobenius norm conservation factor scales as $O(d_{out} d_{in})$ or $O(r^2)$ per expert; and (4) reconstructing the updated task vector $U_{joint} C_k^{(\alpha)} V_{joint}^T$ costs $O(d_{out} d_{in} r)$ FLOPs per expert. Hence, the total complexity of EC-QC-Merge scales as $O(K \cdot d_{out} d_{in} r)$, representing a linear scaling with the number of experts $K$. While more computationally demanding than ThermoMerge, SVD operations are highly parallelizable on modern GPU accelerators, and because merging is a one-time offline process prior to deployment, this modest overhead is highly justified by the massive gains in multi-task representation alignment and performance balance.

\subsection{Limitations and Future Perspectives}
While our physical model merging frameworks are highly promising, we acknowledge several limitations. First, our empirical validation is conducted on a 3-task CIFAR-10 split with a SimpleCNN. Although chosen to isolate weight-space dynamics without architectural confounders, extending our models to large-scale LLMs and ViTs is a critical next step. Second, SVD computations can be intensive for high-dimensional matrices (e.g., $4096 \times 4096$ in LLMs). Exploring randomized or low-rank SVD approximations would mitigate this footprint. Finally, our models introduce thermodynamic parameters ($T, \alpha$); automated tuning or zero-shot hyperparameter selection would enhance practical usability.

\section{Conclusion and Future Work}
\label{sec:conclusion}

In this paper, we introduced \textbf{ThermoMerge}, a unified model merging framework that models singular value spectrum regularization using quantum statistical mechanics. By modeling updates under \textbf{Fermi-Dirac statistics} (FD-Cap), we cap dominant task updates, preventing subspace collapse and ensuring balanced multi-task representation. Furthermore, we demonstrated that warm bosonic statistics (Bose-Einstein Merging at $T=5.0$) can uniquely regularize and align deep multi-layer representation spaces, resolving the Spatial Destruction Fallback.

Future directions include extending ThermoMerge to large language models (LLMs) and Vision Transformers (ViTs), exploring continuous-time thermodynamic flows on Riemannian manifolds, and analyzing the topological invariants of deep weight spaces during parameter fusion.

\newpage
\nocite{*}
\bibliography{submission}
\bibliographystyle{icml2026}

\newpage
\appendix
\onecolumn
\section{Appendix}

\subsection{Proof of Thermodynamic Limits as $T \to \infty$}
We prove that as the thermodynamic temperature $T$ approaches infinity, both the Fermi-Dirac and Bose-Einstein formulations converge to the Isotropic (flat) spectrum, establishing them as continuous generalizations of isotropic merging.

\begin{theorem}
As $T \to \infty$, the unscaled thermodynamic singular values $\tilde{\sigma}_i$ converge to a uniform constant for all states $i$, yielding Isotropic Merging.
\end{theorem}

\begin{proof}
\textbf{Case 1: Fermi-Dirac (FD-Cap)}
The unscaled Fermi-Dirac singular values are given by:
\begin{equation}
    \tilde{\sigma}_i^{FD} = \frac{1}{e^{(-\sigma_i - \mu_{FD})/T} + 1}
\end{equation}
Taking the limit as $T \to \infty$:
\begin{equation}
    \lim_{T \to \infty} \frac{-\sigma_i - \mu_{FD}}{T} = 0
\end{equation}
Therefore, the exponential term approaches $e^0 = 1$:
\begin{equation}
    \lim_{T \to \infty} \tilde{\sigma}_i^{FD} = \frac{1}{1 + 1} = \frac{1}{2} \quad \forall i
\end{equation}
Under Frobenius Norm Conservation, the normalized singular values are scaled by:
\begin{equation}
    \hat{\sigma}_i^{FD} = \tilde{\sigma}_i^{FD} \sqrt{\frac{\sum_{j=1}^r \sigma_j^2}{\sum_{j=1}^r (\tilde{\sigma}_j^{FD})^2}} = \frac{1}{2} \sqrt{\frac{\sum_{j=1}^r \sigma_j^2}{\sum_{j=1}^r (1/4)}} = \frac{1}{2} \sqrt{\frac{\sum_{j=1}^r \sigma_j^2}{r / 4}} = \sqrt{\frac{\sum_{j=1}^r \sigma_j^2}{r}}
\end{equation}
which is exactly a uniform constant for all $i$, matching Isotropic Merging.

\textbf{Case 2: Bose-Einstein}
The unscaled Bose-Einstein singular values are given by:
\begin{equation}
    \tilde{\sigma}_i^{BE} = \frac{1}{e^{(-\sigma_i - \mu_{BE})/T} - 1}
\end{equation}
As $T \to \infty$, we can apply the Taylor expansion of the exponential function $e^x \approx 1 + x$ for small $x$:
\begin{equation}
    e^{(-\sigma_i - \mu_{BE})/T} \approx 1 + \frac{-\sigma_i - \mu_{BE}}{T}
\end{equation}
Substituting this back into the denominator:
\begin{equation}
    \tilde{\sigma}_i^{BE} \approx \frac{1}{\left(1 + \frac{-\sigma_i - \mu_{BE}}{T}\right) - 1} = \frac{T}{-\sigma_i - \mu_{BE}}
\end{equation}
Taking the ratio of any two singular values $i$ and $k$:
\begin{equation}
    \lim_{T \to \infty} \frac{\tilde{\sigma}_i^{BE}}{\tilde{\sigma}_k^{BE}} = \lim_{T \to \infty} \frac{T / (-\sigma_i - \mu_{BE})}{T / (-\sigma_k - \mu_{BE})} = \frac{-\sigma_k - \mu_{BE}}{-\sigma_i - \mu_{BE}}
\end{equation}
Under the condition that the temperature $T$ is much larger than any individual singular value difference, the relative occupancies become approximately uniform. When Frobenius norm normalization is applied, the normalized singular values converge to a flat spectrum:
\begin{equation}
    \hat{\sigma}_i^{BE} \to \sqrt{\frac{\sum_{j=1}^r \sigma_j^2}{r}}
\end{equation}
concluding the proof.
\end{proof}

\subsection{Detailed Hyperparameter Sweeps}
To ensure full reproducibility and provide deep insight into the temperature dynamics of ThermoMerge, we present the complete temperature sweep results ($T \in \{0.1, 1.0, 5.0, 10.0\}$) under the \texttt{classifier\_only} target layer strategy across different merging coefficients $\lambda \in \{0.3, 0.5, 0.7\}$ in \cref{tab:sweep_all}.

\begin{table*}[h]
\caption{Comprehensive temperature sweep results for ThermoMerge and Bose-Einstein Merging under \texttt{classifier\_only} layer strategy. Values represent Task A, B, and C accuracies (\%), Average Accuracy (\%), and Min/Max Task Balance.}
\label{tab:sweep_all}
\vskip 0.15in
\begin{center}
\begin{footnotesize}
\setlength{\tabcolsep}{5pt}
\begin{tabular}{llccccc}
\toprule
$\lambda$ & Method \& Temperature & Task A Acc & Task B Acc & Task C Acc & Avg Acc & Balance \\
\midrule
\textbf{0.3} & ThermoMerge (FD-Suppress, T=0.1) & 24.03\% & 32.63\% & 30.30\% & 28.99\% & 0.74 \\
             & ThermoMerge (FD-Suppress, T=1.0) & 24.93\% & 33.63\% & 29.82\% & 29.46\% & 0.74 \\
             & ThermoMerge (FD-Suppress, T=5.0) & 26.20\% & 36.43\% & 27.65\% & 30.09\% & 0.72 \\
             & ThermoMerge (FD-Suppress, T=10.0) & 26.50\% & 37.07\% & 27.30\% & 30.29\% & 0.71 \\
\cmidrule{2-7}
             & ThermoMerge (FD-Cap, T=0.1) & 35.93\% & 52.20\% & 10.15\% & 32.76\% & 0.19 \\
             & ThermoMerge (FD-Cap, T=1.0) & 30.67\% & 42.97\% & 21.43\% & 31.69\% & 0.50 \\
             & ThermoMerge (FD-Cap, T=5.0) & 27.50\% & 38.70\% & 25.55\% & 30.58\% & 0.66 \\
             & ThermoMerge (FD-Cap, T=10.0) & 27.03\% & 38.13\% & 26.25\% & 30.47\% & 0.69 \\
\cmidrule{2-7}
             & BEMerge (Bose-Einstein, T=0.1) & 39.17\% & 56.70\% & 2.12\% & 32.66\% & 0.04 \\
             & BEMerge (Bose-Einstein, T=1.0) & 37.70\% & 54.83\% & 3.98\% & 32.17\% & 0.07 \\
             & BEMerge (Bose-Einstein, T=5.0) & 35.03\% & 49.23\% & 11.82\% & 32.03\% & 0.24 \\
             & BEMerge (Bose-Einstein, T=10.0) & 34.47\% & 48.43\% & 13.47\% & 32.12\% & 0.28 \\
\midrule
\textbf{0.5} & ThermoMerge (FD-Suppress, T=0.1) & 17.40\% & 20.40\% & 30.77\% & 22.86\% & 0.57 \\
             & ThermoMerge (FD-Suppress, T=1.0) & 20.13\% & 22.83\% & 31.35\% & 24.77\% & 0.64 \\
             & ThermoMerge (FD-Suppress, T=5.0) & 23.47\% & 29.17\% & 30.02\% & 27.55\% & 0.78 \\
             & ThermoMerge (FD-Suppress, T=10.0) & 23.63\% & 30.60\% & 29.98\% & 28.07\% & 0.79 \\
\cmidrule{2-7}
             & ThermoMerge (FD-Cap, T=0.1) & 33.13\% & 55.13\% & 9.72\% & 32.66\% & 0.18 \\
             & ThermoMerge (FD-Cap, T=1.0) & 27.70\% & 41.47\% & 26.00\% & 31.72\% & 0.63 \\
             & ThermoMerge (FD-Cap, T=5.0) & 24.83\% & 33.73\% & 28.93\% & 29.16\% & 0.74 \\
             & ThermoMerge (FD-Cap, T=10.0) & 24.53\% & 32.77\% & 29.52\% & 28.94\% & 0.75 \\
\cmidrule{2-7}
             & BEMerge (Bose-Einstein, T=0.1) & 30.63\% & 60.90\% & 1.20\% & 30.91\% & 0.02 \\
             & BEMerge (Bose-Einstein, T=1.0) & 30.17\% & 60.03\% & 3.25\% & 31.15\% & 0.05 \\
             & BEMerge (Bose-Einstein, T=5.0) & 29.93\% & 53.70\% & 14.75\% & 32.79\% & 0.27 \\
             & BEMerge (Bose-Einstein, T=10.0) & 29.47\% & 52.40\% & 16.82\% & 32.90\% & 0.32 \\
\midrule
\textbf{0.7} & ThermoMerge (FD-Suppress, T=0.1) & 8.30\% & 8.60\% & 26.65\% & 14.52\% & 0.31 \\
             & ThermoMerge (FD-Suppress, T=1.0) & 10.37\% & 11.73\% & 26.65\% & 16.25\% & 0.39 \\
             & ThermoMerge (FD-Suppress, T=5.0) & 14.80\% & 20.50\% & 26.18\% & 20.49\% & 0.57 \\
             & ThermoMerge (FD-Suppress, T=10.0) & 15.93\% & 22.03\% & 26.23\% & 21.40\% & 0.61 \\
\cmidrule{2-7}
             & ThermoMerge (FD-Cap, T=0.1) & 23.37\% & 49.30\% & 1.85\% & 24.84\% & 0.04 \\
             & ThermoMerge (FD-Cap, T=1.0) & 21.40\% & 38.77\% & 22.12\% & 27.43\% & 0.55 \\
             & ThermoMerge (FD-Cap, T=5.0) & 18.60\% & 28.00\% & 25.15\% & 23.92\% & 0.66 \\
             & ThermoMerge (FD-Cap, T=10.0) & 17.67\% & 25.80\% & 25.45\% & 22.97\% & 0.68 \\
\cmidrule{2-7}
             & BEMerge (Bose-Einstein, T=0.1) & 17.73\% & 55.23\% & 0.03\% & 24.33\% & 0.00 \\
             & BEMerge (Bose-Einstein, T=1.0) & 17.57\% & 54.80\% & 0.23\% & 24.20\% & 0.00 \\
             & BEMerge (Bose-Einstein, T=5.0) & 20.67\% & 51.20\% & 9.12\% & 27.00\% & 0.18 \\
             & BEMerge (Bose-Einstein, T=10.0) & 20.20\% & 49.83\% & 11.05\% & 27.03\% & 0.22 \\
\bottomrule
\end{tabular}
\end{footnotesize}
\end{center}
\vskip -0.15in
\end{table*}

\end{document}